\BourbakiExercises{Exercises}

\begin{enumerate}
\def\labelenumi{\arabic{enumi})}
\tightlist
\item
  Let A be a principal ideal domain that is not a field and P a system of representatives consisting of irreducible elements of A. An A-module M is Artinian if and only if it is a torsion module and there exists a finite subset S of P such that the following two conditions are satisfied:
\end{enumerate}

\begin{enumerate}
\def\labelenumi{(\roman{enumi})}
\item
  For every \(\pi \in \mathrm{P} - \mathrm{S}\) , the \(\pi\) -primary component \(\mathrm{M}_{\pi}\) of \(\mathrm{M}\) (VII, §2, No.~2, p.~8) is reduced to 0.
\item
  For every \(\pi\in S\) , the set of elements of M annihilated by \(\pi\) , viewed as a vector space over the field A/ \(\pi\) A, is finite-dimensional.
\end{enumerate}

\begin{enumerate}
\def\labelenumi{\arabic{enumi})}
\setcounter{enumi}{1}
\tightlist
\item
  Let A be a ring, M an A-module of finite length n, and R a set of nilpotent endomorphisms of M that is stable under composition.
\end{enumerate}

\begin{enumerate}
\def\labelenumi{\alph{enumi})}
\item
  If n = 1, then every element of R is zero. If n \textgreater{} 1, prove that there exists a submodule N of M, distinct from 0 and M, that is stable under all elements of R. (We may assume that R has an element \(s \neq 0\) ; choose one such that \(s(M)\) has minimum length, and prove that we have srs = 0 for every \(r \in R\) . Deduce that we can take \(N = s(M)\) if \(Rs = \{0\}\) and \(N = \sum_{r \in R} r(s(M))\) if \(Rs \neq \{0\}\) .) \(^{(1)}\)
\item
  Deduce from a) that there exists a Jordan--Hölder series \((\mathrm{M}_{i})_{0\leqslant i\leqslant n}\) of M such that we have \(r(\mathrm{M}_{i})\subset\mathrm{M}_{i+1}\) for all \(r\in R\) and \(0\leqslant i\leqslant n-1\) . Deduce that we have \(r_{1}\cdots r_{n}=0\) for every sequence \((r_{i})_{1\leqslant i\leqslant n}\) of n elements of R.
\end{enumerate}

\begin{enumerate}
\def\labelenumi{\arabic{enumi})}
\setcounter{enumi}{2}
\tightlist
\item
  \begin{enumerate}
  \def\labelenumii{\alph{enumii})}
  \tightlist
  \item
    Give an example of a commutative field K and an isomorphism \(\varphi\) from K to one of its subfield \(K'\) such that K is finite-dimensional over \(K'\) . On the additive group \(A = K \times K\) , we define a ring structure by setting \((x, y)(x', y') = (xx', xy' + y\varphi(x'))\) (II, §1, p.~382, Exercise 7). The only left ideals of A are \(\{(0, 0)\}\) , \(\{0\} \times K\) , and A, but for every \(K'\) -linear subspace E of K, \(\{0\} \times E\) is a right ideal of A. Deduce that A is left Artinian and left Noetherian but neither right Artinian nor right Noetherian. b) Give an example of a module of finite length whose endomorphism ring is neither left Artinian nor left Noetherian.
  \end{enumerate}
\end{enumerate}

\begin{enumerate}
\def\labelenumi{\alph{enumi})}
\setcounter{enumi}{2}
\tightlist
\item
  Using the same method as in a), give an example of a left and right Artinian ring whose left length is different from its right length.
\end{enumerate}

\begin{enumerate}
\def\labelenumi{\arabic{enumi})}
\setcounter{enumi}{3}
\tightlist
\item
  Let \(\rho : \mathrm{A} \to \mathrm{B}\) be a ring homomorphism.
\end{enumerate}

\begin{enumerate}
\def\labelenumi{\alph{enumi})}
\item
  If, as a right A-module, B is free and nonzero, and if the ring B is left Artinian (resp. left Noetherian), then the ring A is left Artinian (resp. left Noetherian).
\item
  Give an example where A is a commutative field, where B, viewed as a right vector space over A, is finite-dimensional, and where B is neither left Artinian nor left Noetherian (use Exercise 3).
\end{enumerate}

\begin{enumerate}
\def\labelenumi{\arabic{enumi})}
\setcounter{enumi}{4}
\item
  Let n be an integer \(\geqslant 1\) . A ring A is left Artinian (resp. left Noetherian) if and only if the matrix ring \(\mathbf{M}_{n}(\mathrm{A})\) is. If A is left Artinian, then the left length of \(\mathbf{M}_{n}(\mathrm{A})\) is equal to n times that of A.
\item
  Let \(A\) be a ring and \(e\) an idempotent in \(A\) .
\end{enumerate}

\begin{enumerate}
\def\labelenumi{\alph{enumi})}
\item
  Let \(\mathfrak{a}_1\) be a left ideal of the ring \(e\mathrm{A}e\) and \(\mathfrak{a}\) the left ideal of A generated by \(\mathfrak{a}_1\) . We have \(\mathfrak{a}_1 = e\mathfrak{a}e = \mathfrak{a} \cap e\mathrm{A}e\) .
\item
  For every left ideal a of A, we have \(a \cap eAe \subset eae\) ; give an example where this is a strict inclusion (take a matrix ring for A).
\item
  We have \(a \cap eAe = eae\) for every two-sided ideal a of A.
\item
  If the ring A is left Artinian (resp. left Noetherian), then so are the rings eAe and \((1-e)\mathrm{A}(1-e)\) .
\item
  Give an example where eAe and \((1-e)\mathrm{A}(1-e)\) are commutative fields and where A admits infinite strictly increasing sequences and infinite strictly decreasing sequences of two-sided ideals. (Choose an infinite-dimensional vector space V over a commutative field K, endow the additive group \(A=K\times V\times K\) with the ring structure defined by \((\lambda,x,\mu)(\lambda',x',\mu')=(\lambda\lambda',\lambda x'+\mu'x,\mu\mu')\) , and set \(e=(1,0,0)\) .)
\end{enumerate}

\begin{enumerate}
\def\labelenumi{\arabic{enumi})}
\setcounter{enumi}{6}
\tightlist
\item
  \begin{enumerate}
  \def\labelenumii{\alph{enumii})}
  \tightlist
  \item
    Let A be a ring and \(\mathcal{I}\) the set of left ideals of A of the form \(Ae\) , where \(e\) is an idempotent in A. Prove that the following properties are equivalent:
  \end{enumerate}
\end{enumerate}

\begin{enumerate}
\def\labelenumi{(\roman{enumi})}
\item
  Every nonempty subset of \(\mathcal{I}\) has a minimal element.
\item
  Every nonempty subset of \(\mathcal{I}\) has a maximal element.
\item
  There is no infinite family of nonzero ideals belonging to \(\mathcal{I}\) whose sum is direct.
\item
  There is no infinite family \((e_i)\) of nonzero idempotents in A such that \(e_i e_j = 0\) for \(i \neq j\) .
\end{enumerate}

\begin{enumerate}
\def\labelenumi{\alph{enumi})}
\setcounter{enumi}{1}
\item
  Give an example of a ring A in which every nonempty set of monogenous left ideals has a maximal element and a minimal element (so that the equivalent conditions of \(a\) ) are satisfied) even though A is neither left Artinian nor left Noetherian (use Exercise 3).
\item
  Let M be a module. If every nonempty set of finitely generated submodules of M has a maximal element, then M is Noetherian.
\item
  Give an example of a non-Artinian module in which every nonempty set of finitely generated submodules has a minimal element.
\end{enumerate}

\begin{enumerate}
\def\labelenumi{\arabic{enumi})}
\setcounter{enumi}{7}
\item
  Let A be a ring, B and D two subrings of A. Prove that if D is a field and the ring B is left Artinian, \(B \cap D\) is a field.
\item
  Let A be a nonzero commutative ring, and let I be an infinite set. The rings \(\mathrm{A}[(\mathrm{X}_{i})_{i\in\mathrm{I}}]\) and \(\mathrm{A}[[(\mathrm{X}_{i})_{i\in\mathrm{I}}]]\) are neither Artinian nor Noetherian.
\item
  Let A be a ring and a and b two-sided ideals of A. If the rings A/a and A/b are left Artinian (resp. left Noetherian), then so is the ring A/a∩b; by contrast, the ring A/ab is not necessarily left Artinian (resp. left Noetherian).
\item
  Let A, B be rings, M a finitely generated left A-module, P a finitely generated right A-module, and N an (A, B)-bimodule. If the right B-module N is Artinian (resp. Noetherian), then so are the right B-modules \(\operatorname{Hom}_{\mathrm{A}}(\mathrm{M}, \mathrm{N})\) and \(P \otimes_{A} N\) .
\item
  Let \(\mathrm{K}\) be a commutative field. Denote by \((e_n)_{n\in \mathbf{N}}\) the canonical basis of the vector space \(\mathrm{V} = \mathrm{K}^{(\mathbf{N})}\) . We define the sequence \((u_n)_{n\in \mathbf{N}}\) of endomorphisms of \(\mathrm{V}\) by setting
\end{enumerate}

\[
\begin{array}{l l} u _ {0} (e _ {0}) = 0 & \qquad u _ {0} (e _ {1}) = 0 \qquad u _ {0} (e _ {n + 1}) = e _ {n} \text { for } n \geqslant 1 \\ u _ {m} (e _ {0}) = e _ {m} & \qquad u _ {m} (e _ {n}) = 0 \text { for } m \geqslant 1 \text { and } n \geqslant 1. \end{array}
\]

Let \(\mathrm{A}\) be the unital K-subalgebra of \(\operatorname{End}_{\mathrm{K}}(\mathrm{V})\) generated by the sequence \((u_n)_{n \in \mathbf{N}}\) . Prove that the left A-module \(\mathrm{V}\) is Artinian, finitely generated (and even monogenous), but not Noetherian.

\begin{enumerate}
\def\labelenumi{\arabic{enumi})}
\setcounter{enumi}{12}
\tightlist
\item
  \begin{enumerate}
  \def\labelenumii{\alph{enumii})}
  \tightlist
  \item
    Let \(\rho: \mathrm{A} \to \mathrm{B}\) be a ring homomorphism. Assume that A is left Noetherian and that there exists a finite subset S of B whose elements commute with one another and with the elements of \(\rho(\mathrm{A})\) , such that the ring B is generated by \(\rho(\mathrm{A})\) and S. Then B is left Noetherian.
  \end{enumerate}
\end{enumerate}

\begin{enumerate}
\def\labelenumi{\alph{enumi})}
\setcounter{enumi}{1}
\tightlist
\item
  Let A be a nonzero commutative ring. Give an example of an associative unital A-algebra, generated by finitely many elements, that is a neither left nor right Noetherian ring (consider the free associative algebra \(\operatorname{Libas}_{\mathrm{A}}(\mathrm{I})\) (III, §2, No.~7, p.~449, Definition 2), where I is a set of cardinality \(\geqslant 2\) ).
\end{enumerate}

\begin{enumerate}
\def\labelenumi{\arabic{enumi})}
\setcounter{enumi}{13}
\tightlist
\item
  Let V be a vector space of countable infinite dimension over a field K. The set I of endomorphisms of V of finite rank is a two-sided ideal of the ring \(B = \text{End}_{K}(V)\) . Prove that the only two-sided ideals of the ring \(A = B/I\) are \(\{0\}\) and A, but that A is neither left nor right Noetherian.
\end{enumerate}

¶ 15) Let A be a commutative ring and M a finitely generated A-module.

\begin{enumerate}
\def\labelenumi{\alph{enumi})}
\item
  Suppose that for every increasing sequence \((\mathfrak{a}_{n})_{n\in\mathbf{N}}\) of ideals of A, the sequence \((\mathfrak{a}_{n}\mathrm{M})_{n\in\mathbf{N}}\) is stationary. Prove that the A-module M is Noetherian. (Reduce to the case when the A-module M is faithful and when for every ideal \(a\neq0\) of A, the A-module M/aM is Noetherian. There then exists a submodule \(N_{0}\) of M that is maximal among the submodules \(N\subset M\) such that M/N is faithful. Prove that the A-module \(M/N_{0}\) is Noetherian and then that the ring A is Noetherian, and conclude.)
\item
  *Suppose that for every decreasing sequence \((\mathfrak{a}_n)_{n \in \mathbf{N}}\) of ideals of A, the sequence \((\mathfrak{a}_n\mathrm{M})_{n \in \mathbf{N}}\) is stationary. Prove that the A-module M is Artinian. (Reduce to the case when the A-module M is faithful. Prove that A then has only finitely many maximal ideals and, drawing on the proof of Proposition 1 of VIII, p.~173, that the radical of A is nilpotent. Deduce that there exists a finite sequence \((\mathfrak{m}_{1}, \mathfrak{m}_{2}, \cdots, \mathfrak{m}_{n})\) of maximal ideals of A whose product annihilates M. Conclude by induction on \(n.)_{*}\)
\end{enumerate}

\begin{enumerate}
\def\labelenumi{\arabic{enumi})}
\setcounter{enumi}{15}
\tightlist
\item
  \begin{enumerate}
  \def\labelenumii{\alph{enumii})}
  \tightlist
  \item
    Let A be a commutative ring and \(\rho: A \to B\) an injective ring homomorphism that makes B into a finitely generated left A-module. Suppose that for every increasing (resp. decreasing) sequence \((\mathfrak{a}_{n})_{n \in \mathbf{N}}\) of ideals of A, the sequence \((\mathfrak{a}_{n}B)_{n \in \mathbf{N}}\) is stationary (which is the case if the ring B is right Noetherian (resp. right Artinian)). Prove that the ring A is Noetherian (resp. Artinian) (use Exercise 15).
  \end{enumerate}
\end{enumerate}

\begin{enumerate}
\def\labelenumi{\alph{enumi})}
\setcounter{enumi}{1}
\tightlist
\item
  Give an example of a left and right Artinian and Noetherian ring B and a subring A of B such that B is finitely generated as a left and right A-module and that the ring A is neither left nor right Noetherian or Artinian (choose a suitable integral domain C, with field of fractions K, take B to be the matrix ring \(\mathbf{M}_{3}(\mathrm{K})\) and A to be the subring of B consisting of the matrices \((a_{ij})_{1\leqslant i,j\leqslant3}\) with \(a_{ij}=0\) for \(1\leqslant i<j\leqslant3\) and \(a_{22}\in C.\) )
\end{enumerate}

\begin{enumerate}
\def\labelenumi{\arabic{enumi})}
\setcounter{enumi}{16}
\tightlist
\item
  Let E be a monoid, e its identity element, S a subset of E, and \(S'\) the submonoid of E generated by S. We say that S admits a calculus of left fractions on E if the following two (tautological when E is commutative) conditions are satisfied:
\end{enumerate}

\begin{enumerate}
\def\labelenumi{(\roman{enumi})}
\item
  For every \(a \in \mathrm{E}\) and every \(s \in \mathrm{S}\) , there exist \(b \in \mathrm{E}\) and \(t \in \mathrm{S}'\) such that \(ta = bs\) .
\item
  For every \(a \in \mathrm{E}\) , \(b \in \mathrm{E}\) , and \(s \in \mathrm{S}\) such that \(as = bs\) , there exists a \(t \in \mathrm{S}'\) such that \(ta = tb\) .
\end{enumerate}

\(a)\) Suppose that there exists a homomorphism \(u\) from E to a monoid \(\mathrm{E}'\) such that every element of \(u(\mathrm{S})\) is invertible in \(\mathrm{E}'\) and that for every \(a, b\) in E with \(u(a) = u(b)\) , there exists an \(s \in \mathrm{S}'\) such that \(sa = sb\) . Then S admits a calculus of left fractions on E.

\begin{enumerate}
\def\labelenumi{\alph{enumi})}
\setcounter{enumi}{1}
\item
  Suppose that S admits a calculus of left fractions on E. In \(E \times S'\) , we denote by \((a, s) \sim (b, t)\) the relation ``there exist c and d in E such that we have ca = db and \(cs = dt \in S'\) .'' This is an equivalence relation. We set \(S^{-1}E = (E \times S') / \sim\) , and we denote by \(\varepsilon : E \to S^{-1}E\) the mapping that sends \(a \in E\) to the class of \((a, e)\) . There exists a unique monoid structure on \(S^{-1}E\) such that \(\varepsilon\) is a unital homomorphism, that \(\varepsilon(s)\) is invertible for every \(s \in S\) , and that for \((a, s) \in E \times S'\) , the class of \((a, s)\) in \(S^{-1}E\) is equal to \(\varepsilon(s)^{-1}\varepsilon(a)\) . We say that \(S^{-1}E\) is the monoid of left fractions of E associated with S. It coincides with that defined in I, §2, No.~4, p.~18 when the monoid E is commutative.
\item
  Let \(a, b\) be in E. We have \(\varepsilon(a) = \varepsilon(b)\) if and only if there exists an \(s \in S'\) such that \(sa = sb\) . The mapping \(\varepsilon\) is injective (resp. bijective) if and only if every element of S is right regular (resp. invertible).
\item
  Let f be a unital homomorphism from E to a monoid F such that every element of \(f(\mathrm{S})\) is invertible in F. There is a unique homomorphism \(\overline{f}: S^{-1}E \to F\) such that \(f = \overline{f} \circ s\) . If f is injective, then \(\overline{f}\) is injective.
\item
  We say that S admits a calculus of right fractions on E if it admits a calculus of left fractions on the opposite monoid \(E^{o}\) of E. In this case, define the monoid of right fractions \(ES^{-1}\) of E associated with S, observe that we have \((\mathrm{ES}^{-1})^{\circ} = \mathrm{S}^{-1}(\mathrm{E}^{\circ})\) , and rewrite b) for the monoids of right fractions. If S admits a calculus of left and right fractions on E, then the monoids \(S^{-1}E\) and \(ES^{-1}\) are canonically isomorphic.
\end{enumerate}

\begin{enumerate}
\def\labelenumi{\arabic{enumi})}
\setcounter{enumi}{17}
\tightlist
\item
  \begin{enumerate}
  \def\labelenumii{\alph{enumii})}
  \tightlist
  \item
    Let A be a ring, S a subset of A, and S' the smallest subset of A containing 1 and S and stable under multiplication. We say that S admits a calculus of left fractions on A if the following two (tautological when A is commutative) conditions are satisfied:
  \end{enumerate}
\end{enumerate}

\begin{enumerate}
\def\labelenumi{(\roman{enumi})}
\item
  For every \(a \in \mathrm{A}\) and every \(s \in \mathrm{S}\) , there exist \(b \in \mathrm{A}\) and \(t \in \mathrm{S}'\) such that \(ta = bs\) .
\item
  For every \(a \in \mathrm{A}\) and \(s \in \mathrm{S}\) such that \(as = 0\) , there exists a \(t \in \mathrm{S}'\) such that \(ta = 0\) .
\end{enumerate}

This is equivalent to saying that S admits a calculus of left fractions on the monoid obtained by endowing the set A with only multiplication (Exercise 17). Prove that on the multiplicative monoid \(S^{-1}A\) , there then exists a unique addition such that \(S^{-1}A\) , endowed with this addition and its multiplication, is a ring and that the canonical mapping \(\varepsilon: A \to S^{-1}A\) (loc. cit.) is a ring homomorphism. We say that \(S^{-1}A\) is the ring of left fractions of A associated with S. When A is commutative, this coincides with the ring defined in I, §8, No.~12, p.~113.

\begin{enumerate}
\def\labelenumi{\alph{enumi})}
\setcounter{enumi}{1}
\item
  The set S admits a calculus of left fractions on A if and only if there exist a ring B and a ring homomorphism \(u: A \to B\) such that every element of \(u(S)\) is invertible in B and that for every \(a \in u^{-1}(0)\) , there exists an \(s \in S\) such that sa = 0.
\item
  Suppose that S admits a calculus of left fractions on A. The kernel of the canonical homomorphism \(\varepsilon : \mathrm{A} \to \mathrm{S}^{-1}\mathrm{A}\) is the set of \(a \in \mathrm{A}\) for which there exists an \(s \in \mathrm{S}\) with \(sa = 0\) .
\end{enumerate}

Let \(f\) be a homomorphism from A to a ring B such that every element of \(f(\mathrm{S})\) is invertible in B. There exists a unique homomorphism \(\overline{f} : \mathrm{S}^{-1}\mathrm{A} \to \mathrm{F}\) such that \(f = \overline{f} \circ s\) . If \(f\) is injective, then \(\overline{f}\) is injective.

\begin{enumerate}
\def\labelenumi{\alph{enumi})}
\setcounter{enumi}{3}
\item
  Formulate the statements analogous to a), b), c) for a calculus of right fractions; if S admits a calculus of right fractions on A, we denote the ring of right fractions of A by AS \(^{-1}\) . In this case, we have \((\mathrm{AS}^{-1})^{\circ} = \mathrm{S}^{-1}(\mathrm{A}^{\circ})\) . If S admits a calculus of left and right fractions on A, then the rings S \(^{-1}\) A and AS \(^{-1}\) are canonically isomorphic.
\item
  Let A be a nonzero ring without zero divisors. Then A admits a field of left fractions (I, §9, p.~177, Exercise 15) if and only if the set S of nonzero elements of A admits a calculus of left fractions on A, and in this case, \(S^{-1}A\) is the field of left fractions of A.
\end{enumerate}

\begin{enumerate}
\def\labelenumi{\arabic{enumi})}
\setcounter{enumi}{18}
\tightlist
\item
  Let A be a ring, S a subset of A that admits a calculus of left fractions on A (Exercise 18), and S' the smallest subset of A containing 1 and S and stable under multiplication. We denote by \(S^{-1}A\) the ring of left fractions of A associated with S and by \(\varepsilon:A\to S^{-1}A\) the canonical homomorphism.
\end{enumerate}

\begin{enumerate}
\def\labelenumi{\alph{enumi})}
\item
  Let M be an A-module. We denote the \(S^{-1}A\) -module \(S^{-1}A \otimes_{A} M\) by \(S^{-1}M\) , and we call it the module of left fractions of M associated with S. Prove that every element of \(S^{-1}M\) can be written as \(\varepsilon(s)^{-1} \otimes x\) with \(x \in M\) and \(s \in S'\) , and that for x, y in M and s, t in \(S'\) , the relation \(\varepsilon(s)^{-1} \otimes x = \varepsilon(t)^{-1} \otimes y\) is equivalent to the existence of a pair \((c, d)\) of elements of A such that we have cx = dy and \(cs = dt \in S'\) .
\item
  We denote by \(\varepsilon_{M}\) the canonical A-linear mapping \(x \mapsto 1 \otimes x\) from M to \(S^{-1}M\) . Its image generates the \(S^{-1}A\) -module \(S^{-1}M\) ; its kernel consists of the elements \(x \in M\) for which there exists an \(s \in S'\) such that sx = 0. The mapping \(\varepsilon_{M}\) is injective (resp. bijective) if and only if for every \(s \in S\) , the homothety \(x \mapsto sx\) of M is injective (resp. bijective).
\end{enumerate}

\(c)\) We have \(\mathrm{S}^{-1}\mathrm{M} = 0\) if and only if for every \(x\in \mathrm{M}\) , there exists an \(s\in \mathrm{S}\) such that \(sx = 0\) .

\begin{enumerate}
\def\labelenumi{\alph{enumi})}
\setcounter{enumi}{3}
\item
  Let N be a submodule of M. The canonical \(S^{-1}A\) -linear mapping from \(S^{-1}N\) to \(S^{-1}M\) is injective. (*In other words, the right A-module \(S^{-1}A\) is flat.*) We identify \(S^{-1}N\) with its image in \(S^{-1}M\) under this mapping. The submodule \(\bar{\varepsilon}_{\mathrm{M}}^{-1}(\mathrm{S}^{-1}\mathrm{N})\) of M consists of the elements \(x \in M\) for which there exists an \(s \in S'\) with \(sx \in N\) . We say that it is the left saturation of N for S.
\item
  The mapping \(N^{\prime} \mapsto \varepsilon_{M}^{-1}(N^{\prime})\) is an isomorphism from the ordered set of \(S^{-1}A\) -submodules of \(S^{-1}M\) onto the ordered set of A-submodules of M that are left saturated, (i.e., equal to their saturation) for S. The inverse isomorphism is \(N \mapsto S^{-1}N\) . f) If the A-module M is Noetherian, then the \(S^{-1}A\) -module \(S^{-1}M\) is Noetherian. If the A-module M is Artinian, then the \(S^{-1}A\) -module \(S^{-1}M\) is Artinian and the mapping \(\varepsilon_{M}: M \to S^{-1}M\) is surjective.
\end{enumerate}

\(g\) ) If the ring A is left Noetherian (resp. left Artinian), then so is the ring \(\mathrm{S}^{-1}\mathrm{A}\) . \(h\) ) Formulate the statements analogous to those of this exercise for the modules of right fractions of right modules.

\begin{enumerate}
\def\labelenumi{\arabic{enumi})}
\setcounter{enumi}{19}
\item
  Prove that a left Noetherian ring A without zero divisors admits a field of left fractions (Exercise 18) (if a and s are nonzero elements of A, use the fact that the sequence of ideals \((As + Asa + \cdots + Asa^{n})_{n \in \mathbf{N}}\) is stationary to prove that \(Aa \cap As\) is not reduced to 0).
\item
  Let A be a ring, \(\sigma\) an endomorphism of the ring A, and \(d\) a nilpotent endomorphism of the additive group of A satisfying the relation \(d(ab) = \sigma(a)d(b) + d(a)b\) . \(a)\) Show that the set \(S = \{X\}\) admits a calculus of left fractions (Exercise 18) on the ring \(B = A[X]_{\sigma,d}\) (VIII, p.~11). If \(\sigma\) is injective, then the canonical homomorphism from B to \(S^{-1}B\) is injective and we identify B with a subring of \(S^{-1}B\) . If \(\sigma\) is bijective, then S also admits a calculus of right fractions on B, and every element of \(S^{-1}B\) can be written uniquely as \(\sum_{n\in\mathbf{Z}}a_nX^n\) , where \((a_n)_{n\in\mathbf{Z}}\) is a family of elements of A with finite support. For \(a \in \mathrm{A}\) , write the element \(\mathrm{X}^{-1}a\) of \(\mathrm{S}^{-1}\mathrm{B}\) in this form.
\end{enumerate}

\begin{enumerate}
\def\labelenumi{\alph{enumi})}
\setcounter{enumi}{1}
\item
  Prove that the multiplication in the ring \(A[X]_{\sigma,d}\) extends by continuity to a multiplication on the set \(A[[X]]\) of formal power series in X with coefficients in A. The additive group \(A[[X]]\) endowed with this multiplication is a ring that we denote by \(A[[X]]_{\sigma,d}\) (or simply \(A[[X]]_{\sigma}\) if d=0). Give an example where a formal power series \(u\in A[[X]]\) with constant term equal to 1 has neither a left nor a right inverse in the ring \(A[[X]]_{\sigma,d}\) .
\item
  Show that the invertible elements of \(A[[X]]_{\sigma}\) are those whose constant term is invertible in A.
\item
  Formulate and prove the analog of a) after replacing the ring \(A[X]_{\sigma,d}\) with the ring \(A[[X]]_{\sigma,d}\) .
\item
  Suppose that A is a field and \(\sigma\) an automorphism of A. Then \(\mathrm{S}^{-1}\mathrm{A}[[\mathrm{X}]]_{\sigma}\) (where \(\mathrm{S} = \{\mathrm{X}\}\) ) is a field, and each of its elements can be written uniquely as \(\sum_{n\in \mathbf{Z}}a_n\mathrm{X}^n\) , where \((a_{n})_{n\in \mathbf{Z}}\) is a family of elements of A having only finitely many nonzero terms with index \(< 0\) . We denote this field by \(\mathrm{A}((\mathrm{X}))_{\sigma}\) .
\end{enumerate}

\(f\) ) Suppose that A is a commutative field. Describe the center of the field \(\mathrm{A}((\mathrm{X}))_{\sigma}\) . Deduce an example of a field of infinite degree over its center.

\begin{enumerate}
\def\labelenumi{\arabic{enumi})}
\setcounter{enumi}{21}
\tightlist
\item
  Let G be a group and H a normal subgroup of G such that the quotient group G/H is isomorphic to Z. Let g be an element of G whose canonical image in G/H generates G/H. Let C be a commutative ring, A the algebra C{[}H{]} of H over C (III, §2, No.~6, p.~446), \((e_{h})_{h\in\mathrm{H}}\) its canonical basis, and \(\sigma\) the automorphism of A that sends \(e_{h}\) to \(e_{ghg^{-1}}\) .
\end{enumerate}

\begin{enumerate}
\def\labelenumi{\alph{enumi})}
\item
  Prove that there exists a unique homomorphism \(u\) from the ring \(\mathrm{B} = \mathrm{A}[\mathrm{X}]_{\sigma}\) to the ring \(\mathrm{C}[\mathrm{G}]\) that extends the canonical injection of \(\mathrm{A} = \mathrm{C}[\mathrm{H}]\) into \(\mathrm{C}[\mathrm{G}]\) and sends \(\mathrm{X}\) to the element \(e_g\) of the canonical basis of \(\mathrm{C}[\mathrm{G}]\) .
\item
  Set \(S = \{X\}\) . Prove that u extends to an isomorphism from the ring \(S^{-1}B\) (cf.~Exercise 21) to C{[}G{]}. Deduce that if the ring C{[}H{]} is left Noetherian, then so is the ring C{[}G{]}.
\end{enumerate}

\begin{enumerate}
\def\labelenumi{\arabic{enumi})}
\setcounter{enumi}{22}
\tightlist
\item
  Let A be a ring, I a set, \(\symbf{\sigma} = (\sigma_{i})_{i \in \mathrm{I}}\) a family of endomorphisms of the ring A, and \(\symbf{d} = (d_{i})_{i \in \mathrm{I}}\) a family of endomorphisms of the additive group of A, satisfying the relations
\end{enumerate}

\[
\sigma_ {i} \circ \sigma_ {j} = \sigma_ {j} \circ \sigma_ {i} \quad \sigma_ {i} \circ d _ {j} = d _ {j} \circ \sigma_ {i} \quad d _ {i} \circ d _ {j} = d _ {j} \circ d _ {i}
\]

for \(i\in \mathrm{I},j\in \mathrm{I},i\neq j\) and

\[
d _ {i} (a b) = \sigma_ {i} (a) d _ {i} (b) + d _ {i} (a) b
\]

for \(i \in I\) , \(a \in A\) , \(b \in A\) .

\begin{enumerate}
\def\labelenumi{\alph{enumi})}
\tightlist
\item
  Prove that there exists a unique ring B whose additive group is the additive group A{[}X{]} of polynomials with coefficients in A in the family of variables \(\mathbf{X} = (\mathrm{X}_{i})_{i \in \mathrm{I}}\) and whose product has the following properties:
\end{enumerate}

\begin{enumerate}
\def\labelenumi{(\roman{enumi})}
\item
  The elements \(\mathrm{X}_i\) ( \(i \in \mathrm{I}\) ) of B are pairwise permutable.
\item
  For every \(a \in A\) and every multi-index \(\boldsymbol{\nu} = (\nu_i)_{i \in I} \in \mathbf{N}^{(I)}\) , the product in B, in this order, of a and a finite sequence of \(\nu_i\) elements equal to \(X_i\) for each \(i \in I\) is the polynomial \(aX^\nu\) .
\item
  For every \(a \in \mathrm{A}\) and every \(i \in \mathrm{I}\) , we have the relation \(\mathrm{X}_i a = \sigma_i(a)\mathrm{X}_i + d_i(a)\) in the ring B.
\end{enumerate}

\begin{enumerate}
\def\labelenumi{\alph{enumi})}
\setcounter{enumi}{1}
\item
  The ring B is denoted by \(A[X]_{\sigma,d}\) . It has the following universal property: given a ring C, a ring homomorphism \(f: A \to C\) , and a family \((x_i)_{i \in I}\) of elements of C that are pairwise permutable such that \(x_i f(a) = f(\sigma_i(a)) x_i + f(d_i(a))\) for every \(a \in A\) and every \(i \in I\) , there exists a unique ring homomorphism \(g: B \to C\) extending f such that \(g(X_i) = x_i\) for \(i \in I\) .
\item
  Let J be a subset of I. Denote the ring \(\mathrm{A}[(\mathrm{X}_{i})_{i\in\mathrm{J}}]_{(\sigma_{i})_{i\in\mathrm{J}},(d_{i})_{i\in\mathrm{J}}}\) by \(B'\) . For every \(i\in I-J\) , there exist a unique endomorphism \(\sigma_{i}'\) of the ring \(B'\) and a unique endomorphism \(d_{i}'\) of the additive group of \(B'\) such that
\end{enumerate}

\[
\begin{array}{l l} \sigma_ {i} ^ {\prime} (a) = \sigma_ {i} (a) & d _ {i} ^ {\prime} (a) = d _ {i} (a) \qquad \text {for a\ensuremath{\in} A}, \\ \sigma_ {i} ^ {\prime} (\mathrm{X} _ {j}) = \mathrm{X} _ {j} & d _ {i} ^ {\prime} (\mathrm{X} _ {j}) = 0 \qquad \text {for j\ensuremath{\in} J}, \\ d _ {i} ^ {\prime} (\mathrm{PQ}) = \sigma_ {i} ^ {\prime} (\mathrm{P}) d _ {i} ^ {\prime} (\mathrm{Q}) + d _ {i} ^ {\prime} (\mathrm{P}) \mathrm{Q} & \text {for P\ensuremath{\in} B , Q\ensuremath{\in} B}. \end{array}
\]

Set \(\sigma' = (\sigma_i')_{i \in \mathbf{I} - \mathbf{J}}\) and \(d' = (d_i')_{i \in \mathbf{I} - \mathbf{J}}\) . The canonical bijection from \(\mathrm{A}[(\mathrm{X}_i)_{i \in \mathbf{I}}]\) to \(\mathrm{A}[(\mathrm{X}_j)_{j \in \mathbf{J}}][(\mathrm{X}_i)_{i \in \mathbf{I} - \mathbf{J}}]\) is an isomorphism from the ring \(\mathrm{B} = \mathrm{A}[(\mathrm{X}_i)_{i \in \mathbf{I}}]_{\sigma, d}\) to the ring \(\mathrm{B}'[(\mathrm{X}_i)_{i \in \mathbf{I} - \mathbf{J}}]_{\sigma', d'}\) .

\(d)\) If the ring A is left Noetherian, then the set I is finite, and if for every \(i \in \mathrm{I}\) , the endomorphism \(\sigma_i\) of A is bijective, then the ring \(\mathrm{B} = \mathrm{A}[\mathbf{X}]_{\sigma,d}\) is left Noetherian. \(e)\) Suppose that for every \(i \in \mathrm{I}\) , the morphism \(\sigma_i\) is the identity mapping on A. Let E be the endomorphism ring of the additive group of A. Show that there exists a unique ring homomorphism \(\psi\) from \(\mathrm{B} = \mathrm{A}[\mathbf{X}]_{\sigma,d}\) to E such that \(\psi(\mathrm{X}_i) = d_i\) for \(i \in \mathrm{I}\) and that for all \(a \in \mathrm{A}\) , the morphism \(\psi(a)\) is the left homothety \(b \mapsto ab\) of A. \(f)\) Keep the assumptions of \(e)\) . Suppose, moreover, that A is a polynomial ring \(\mathrm{C}[(\mathrm{T}_i)_{i \in \mathrm{I}}]\) with coefficients in a (not necessarily commutative) ring C and that for every \(i \in \mathrm{I}\) , the morphism \(d_i\) is the derivation \(\mathrm{P} \mapsto \frac{\partial\mathrm{P}}{\partial\mathrm{T}_i}\) of A. Prove that if C is nonzero and without torsion on \(\mathbf{Z}\) , the homomorphism \(\psi: \mathrm{B} \to \mathrm{E}\) is injective. Prove that if C is a ring of characteristic \(p > 0\) (V, §1, No.~2, p.~2), the kernel of \(\psi\) is the left ideal of B generated by the element \(\mathrm{X}_i^p\) ( \(i \in \mathrm{I}\) ).

¶ 24) Let G be a group, e its identity element, and A a nonzero commutative ring. Denote the algebra of the group G over A by A{[}G{]} (III, §2, No.~6, p.~446).

\begin{enumerate}
\def\labelenumi{\alph{enumi})}
\item
  Let I be a finite set, \((\mathrm{G}_{i})_{i\in\mathrm{I}}\) a family of subgroups of G, and \((g_{i})_{i\in\mathrm{I}}\) a family of elements of G. Prove that if G is the union of the family \((g_{i}\mathrm{G}_{i})_{i\in\mathrm{I}}\) , one of the subgroups \(G_{i}\) has finite index in G.
\item
  Suppose that every conjugacy class in G distinct from \(\{e\}\) is infinite and that the ring A is an integral domain. Prove that if a and b are two nonzero two-sided ideals of the ring A{[}G{]}, we have \(ab \neq 0\) . (Choose elements \(a = \sum a_g g\) and \(b = \sum b_g g\) in a and b with \(a_e \neq 0\) and \(b_e \neq 0\) . Using a), prove that, if necessary replacing a with \(hah^{-1}\) for suitable \(h \in G\) , we may assume that \(a_g b_{g^{-1}} = 0\) for every \(g \neq e\) in G.) *c) Prove that the ring A{[}G{]} is left (resp. right) Artinian if and only if the ring A is Artinian and the group G is finite. (If A{[}G{]} is left Artinian, deduce from Theorem 1 of VIII, p.~6 that the group G has finite length; to prove that it is finite, reduce to the case when A is a field and G a simple group, and then use b) and Proposition 1 of VIII, p.~173).*
\end{enumerate}

\begin{enumerate}
\def\labelenumi{\arabic{enumi})}
\setcounter{enumi}{24}
\tightlist
\item
  Keep the notation of Exercise 24.
\end{enumerate}

\begin{enumerate}
\def\labelenumi{\alph{enumi})}
\item
  For any subgroup H of G, let \(I_{H}\) be the kernel of the canonical surjection \(\mathbf{Z}[G] \to \mathbf{Z}^{(G/H)}\) . Then the mapping \(H \mapsto I_{H}\) is the increasing injection from the set of subgroups of G to the set of the left ideals of \(Z[G]\) ; if H is a normal subgroup of G, the ideal \(I_{H}\) is two-sided.
\item
  Suppose that the ring A{[}G{]} is left Noetherian. It is then also right Noetherian. For every subgroup H of G, the ring A{[}H{]} is left Noetherian; in particular, the ring A is Noetherian. If H is normal in G, then the ring A{[}G/H{]} is left Noetherian. Every subgroup of G admits a finite generating system.
\item
  Prove that the following properties are equivalent:
\end{enumerate}

\begin{enumerate}
\def\labelenumi{(\roman{enumi})}
\item
  The group G admits a composition series whose quotients are finite or isomorphic to Z.
\item
  The group G admits a composition series \((\mathrm{G}_i)_{0\leqslant i\leqslant n}\) such that \(\mathrm{G}_1\) has finite index in G and \(\mathrm{G}_i / \mathrm{G}_{i + 1}\) is isomorphic to \(\mathbf{Z}\) for \(0\leqslant i\leqslant n - 1\) .
\end{enumerate}

When they are satisfied and the ring A is Noetherian, the ring A{[}G{]} is left Noetherian (use Exercise 22).

\begin{enumerate}
\def\labelenumi{\arabic{enumi})}
\setcounter{enumi}{25}
\item
  Let K be a commutative field, A the polynomial ring K{[}T{]}, and \(\sigma\) the endomorphism \(\mathrm{P(T)}\mapsto\mathrm{P(T^{2})}\) of the ring A. The ring \(A[X]_{\sigma}\) is not left Noetherian even though the ring A is Noetherian.
\item
  Let D be a field, \(\sigma\) an automorphism of D, \(d\) an endomorphism of the additive group of D satisfying relation (1) of VIII, p.~9, and A the ring \(\mathrm{D}[\mathrm{X}]_{\sigma,d}\) (VIII, p.~11). Given a nonzero element \(\mathrm{P} = \sum_{n}a_{n}\mathrm{X}^{n}\) of A, the greatest integer \(n\) such that \(a_{n}\neq 0\) is called the degree of P and denoted by \(\deg (\mathrm{P})\) . We set \(\deg (0) = -\infty\) .
\end{enumerate}

\begin{enumerate}
\def\labelenumi{\alph{enumi})}
\item
  Let F, G be elements of A with \(G \neq 0\) ; show that there exist well-determined elements Q, R (resp. \(Q', R'\) ) of A such that \(F = QG + R\) and \(\deg(R) < \deg(G)\) (resp. \(F = GQ' + R'\) and \(\deg(R') < \deg(G)\) ).
\item
  Deduce that every left (resp. right) ideal of A is monogenous. If F, G, H, K are elements of A satisfying AF + AG = AH and AF ∩ AG = AK, then we have \(\deg(\mathrm{F}) + \deg(\mathrm{G}) = \deg(\mathrm{H}) + \deg(\mathrm{K})\) (observe that the D-vector space A/AF has dimension \(\deg(\mathrm{F})\) ).
\item
  The ideal AF of A is maximal if and only if F is irreducible, that is, is not the product of two nonconstant polynomials.
\end{enumerate}

\begin{enumerate}
\def\labelenumi{\arabic{enumi})}
\setcounter{enumi}{27}
\tightlist
\item
  Let B be a ring, A a subring of B, and x an element of B such that we have \(A + Ax = A + xA\) and that the ring B is generated by \(A \cup \{x\}\) .
\end{enumerate}

\begin{enumerate}
\def\labelenumi{\alph{enumi})}
\item
  Let M be a Noetherian left A-module; prove that the B-module \(\mathrm{M}_{(\mathrm{B})} = \mathrm{B} \otimes_{\mathrm{A}} \mathrm{M}\) is Noetherian. In particular, if the ring A is left Noetherian, then so is B.
\item
  Suppose that the A-module M has finite length n; prove that every B-submodule of \(\mathrm{M}_{(\mathrm{B})}\) admits a generating subset of cardinality \(\leqslant n\) (adapt the proof of Theorem 1 of VIII, p.~6).
\end{enumerate}

